\scanpage{99}
% archive foliation: 100
- 48 -\hfill n\textsuperscript{o} 380

sépare les points de $\Omega$.

PROPOSITION 1.- \uline{Si} $\Omega$ \uline{est complètement régulier, l'application canonique de} $\Omega$ \uline{dans le spectre de} $\mathcal B(\Omega)$ \uline{est un homéomorphisme de} $\Omega$ \uline{sur un sous-espace dense de} $\operatorname{Sp}(\mathcal B(\Omega))$.

En effet, si $\Omega$ est complètement régulier, alors $\mathcal B(\Omega)$ sépare les points de $\Omega$, donc l'application canonique est injective. La topologie de $\Omega$ étant la topologie la \quad fine pour laquelle les éléments de $\mathcal B(\Omega)$ sont continus, l'application canonique est un homéomorphisme de $\Omega$ sur son image dans $\operatorname{Sp}(\mathcal B(\Omega))$. Montrons maintenant que cette image $\Omega'$ est dense. Supposons que $\eta\in\operatorname{Sp}(\mathcal B(\Omega))$ ne soit pas adhérent à $\Omega'$. Il existe alors une famille finie $(f_i)_{i\in I}$ d'éléments de $\mathcal B(\Omega)$ et un nombre $\varepsilon>0$ tels que pour tout $\xi\in\Omega'$, il existe un indice $j\in I$ pour lequel $|f_j(\xi)-f_j(\eta)|>\varepsilon$. On peut supposer la famille $(f_i)_{i\in I}$ choisie de sorte que, l'un au moins des $a_i=\widehat f_i(\eta)$ soit $\ne0$. Posons $b=\sum_{i\in I}a_i\bar a_i$ et
\[
u=b^{-1}\sum_i(a_i-f_i)(\bar a_i-\bar f_i).
\]

Pour tout $t\in\Omega$, $\xi_t\in\Omega'$ et par suite
\[
\begin{aligned}
u(t)=\xi_t(u)&=b^{-1}\sum_i(a_i-f_i(t))(\bar a_i-\overline{f_i(t)})\\
&=b^{-1}\sum_i(\widehat f_i(\eta)-\widehat f_i(\xi_t))(\overline{\widehat f_i(\eta)-\widehat f_i(\xi_t)})\geq b^{-1}\varepsilon^2.
\end{aligned}
\]
Il en résulte que $u$ est inversible dans $\mathcal B(\Omega)$. Par suite $\widehat u$ doit être inversible dans $(\operatorname{Sp}(\mathcal B(\Omega)))$. Or ceci est impossible puisque $\widehat u(\eta)=\eta(u)=0$.

Lorsque $\Omega$ est un espace complètement régulier, le spectre de l'algèbre des fonctions continues bornées sur $\Omega$ est appelé le compactifié modèle Stone Cech de $\Omega$.
(cf. Top. Gén., Chap.9, \S\ 1, Exercices 7 à 13).

COROLLAIRE.- \uline{Si} $\Omega$ \uline{est un espace compact, l'application canonique de} $\Omega$ \uline{dans le spectre de} $\mathcal B(\Omega)$ \uline{est un homéomorphisme}.

En effet, le spectre de $\mathcal B(\Omega)$ est séparé.

\(-:-:-:-\)
